\subsection{Questão 51}

\subsubsection{Enunciado}

Um anel de Rowland é feito de um material ferromagnético. O anel tem seção reta circular, com um raio interno de \(5{,}0\,\text{cm}\) e um raio externo de \(6{,}0\,\text{cm}\), e uma bobina primária enrolada no anel possui \(400\) espiras.

\begin{enumerate}
    \item[(a)] Qual deve ser a corrente na bobina para que o módulo do campo toróide tenha o valor
    \[
    B_0 = 0{,}20\,\text{mT}?
    \]

    \item[(b)] Uma bobina secundária enrolada no anel possui \(50\) espiras e uma resistência de \(8{,}0\,\Omega\). Se para esse valor de \(B_0\) temos
    \[
    B_M = 800B_0,
    \]
    qual é o valor da carga que atravessa a bobina secundária quando a corrente na bobina primária começa a circular?
\end{enumerate}


\subsubsection{Resolução}


\vspace{1cm}
